% SEGMENT OLP-0596-S01
\documentclass[../../../include/open-logic-section]{subfiles}

\begin{document}

\olfileid{sth}{choice}{woproblem}
\olsection{நன்கு வரிசைப்படுத்தல் சிக்கல்}

நன்கு வரிசைப்படுத்தலை ஏற்றுக்கொள்கிறோமா என்பதிலேயே
பல முடிவுகள் தங்கியுள்ளன. ஆனால் இந்தக் கொள்கை பற்றிய
விவாதம் பெரும்பாலும் அதற்குச் சமமான தேர்வின் அடிகோளத்தை
மையமாகக் கொண்டுள்ளது. ஆகவே இப்போது அதைப் பார்ப்போம்;
இரண்டும் சமம் என்பதையும் நிறுவுவோம்.

\citeyear{Cantor1883} இல், காண்டோர் நன்கு வரிசைப்படுத்தல்
அடிகோளத்தை ஆதரித்து, அது தனக்கு அடிப்படையானதாகவும்,
விளைவுகளில் வளமானதாகவும், அதன் பொதுச் செல்லுபடியில்
குறிப்பாகக் கவனிக்கத்தக்கதாகவும் தோன்றும் ``சிந்தனை விதி''
என்று அழைத்தார் (\citeauthor{Potter2004}
\citeyear[p.~243]{Potter2004} இல் மேற்கோள் காட்டப்பட்டுள்ளது).
ஆனால் பின்னர் இந்த ``அடிகோளம்'' நிறுவப்பட வேண்டும் என்று
காண்டோரே கருதினார். கணிதச் சமூகமும் அவ்வாறே கருதியது.

இந்தச் சிக்கலை \citeyear{Zermelo1904} இல் செர்மேலோ
``தீர்த்தார்''. அவருடைய தீர்வை விளக்கச் சில வரையறைகள் தேவை.

\begin{defn}
எல்லா $x \in \dom{f}$ க்கும் $f(x) \in x$ எனில்,
சார்பு $f$ ஒரு \emph{தேர்வுச் சார்பு}. மேலும் $f$ ஒரு தேர்வுச்
சார்பாகவும் $\dom{f} = A \setminus \{\emptyset\}$ ஆகவும்
இருந்தால், $f$ ஐ $A$ க்கான \emph{தேர்வுச் சார்பு} என்கிறோம்.
\end{defn}

உள்ளுணர்வாக, ஒவ்வொரு காலியற்ற கணம்
$x \in A$ இலிருந்தும் $A$ க்கான தேர்வுச் சார்பு $x$ இன்
ஒரு குறிப்பிட்ட உறுப்பான $f(x)$ ஐத் \emph{தேர்ந்தெடுக்கிறது}. அப்போது
தேர்வின் அடிகோளம்:

\begin{axiom}[தேர்வு]
	ஒவ்வொரு கணத்துக்கும் ஒரு தேர்வுச் சார்பு உள்ளது.
\end{axiom}

% SEGMENT OLP-0596-S02
தேர்வு நன்கு வரிசைப்படுத்தலைத் தரும் என்றும், மறுதிசையும்
உண்மை என்றும் செர்மேலோ காட்டினார்:

\begin{thm}[$\ZF$ இல்]\ollabel{thmwochoice}
நன்கு வரிசைப்படுத்தலும் தேர்வும் சமமானவை.
\end{thm}

\begin{proof}
\emph{இடமிருந்து வலம்.} $A$ கணங்களின் கணம் என்க.
ஒன்றிப்பு அடிகோளத்தால் $\bigcup A$ உள்ளது. ஆகவே நன்கு
வரிசைப்படுத்தலால் $\bigcup A$ ஐ நன்கு வரிசைப்படுத்தும்
$<$ என்ற வரிசை உள்ளது. இப்போது
$f(x) = \text{$<$-மீச்சிறிய உறுப்பு }x$ என்க.
இது $A$ க்கான தேர்வுச் சார்பு.

% SEGMENT OLP-0596-S03
\emph{வலமிருந்து இடம்.} $A$ ஐ நிலைப்படுத்துக.
தேர்வின்படி $\Pow{A} \setminus \{\emptyset\}$ க்கான
ஒரு தேர்வுச் சார்பு $f$ உள்ளது.
% SOURCE CORRECTION TA-STH-042: A காலியாக இருந்தால் தனியே கையாள்க; இல்லையேல் f(A) வரையறுக்கப்படாது.
காலிக் கணத்தை நன்கு வரிசைப்படுத்துவது
உடனடியாகும். இனி காலியற்ற கணத்தைக் கருதுவோம்.
முடிவிலிக்கடந்த மீள்வரையறையைப் பயன்படுத்திப் பின்வரும்
சார்பை வரையறுக்கவும்:
\begin{align*}
	g(0) &= f(A)\\
	g(\alpha) &=
		\begin{cases}
			\text{நிறுத்து!{}} &\text{எனில் }A = \funimage{g}{\alpha}\\
			f(A \setminus \funimage{g}{\alpha}) & \text{மற்றபடி}\\
		\end{cases}
\end{align*}
``நிறுத்து!'' என்பது விரிவான மீள்வரையறைக்கு மாற்றான
சுருக்கக் குறிப்பு மட்டுமே. அதாவது, முதன்முதலாக
$A = \funimage{g}{\alpha}$ ஆகும் நிலையில்,
% SOURCE CORRECTION TA-STH-043: நிறுத்தத்துக்குப் பிந்தைய குறிகளுக்கு மட்டுமே நிரப்பு மதிப்பை இடுக.
எல்லா $\alpha \leq \delta$ க்கும் $g(\delta) = A$ என
வரையறுக்கலாம்; முன்னைய மதிப்புகளை மாற்றக் கூடாது.
முதலில் இது ஒரு \emph{பதத்தை} மட்டுமே வரையறுக்கிறது
என்று உறுதியாகச் சொல்ல முடியும்
(\olref[sth][spine][recursion]{transrecursionschema} க்குப்
பிந்தைய குறிப்புகளைப் பார்க்கவும்). ஆனால் வேண்டிய
சார்பு இருப்பதைக் கீழே காட்டுவோம்.

% SEGMENT OLP-0596-S04
$f$ ஒரு தேர்வுச் சார்பு என்பதால், நிறுத்தத்துக்கு முன்
வரையறுக்கப்படும் ஒவ்வொரு $\alpha$ க்கும்
$g(\alpha) = f(A \setminus \funimage{g}{\alpha}) \in A \setminus
\funimage{g}{\alpha}$; அதாவது
$g(\alpha) \notin \funimage{g}{\alpha}$.
எனவே $g(\alpha) = g(\beta)$ எனில்,
$g(\beta) \notin \funimage{g}{\alpha}$; அதாவது
$\beta \notin \alpha$. அதேபோல் $\alpha \notin \beta$.
மூவகைப் பிரிவின்படி $\alpha = \beta$. ஆகவே
% SOURCE CORRECTION TA-STH-044: ஒருமுகத்தன்மை நிறுத்துமுன் உள்ள ஆரம்பத் துண்டுக்கே உரியது.
நிறுத்துமுன் கிடைக்கும் $g$ இன் தொடக்கத் துண்டு
!!{injective} ஆகும்.

% SEGMENT OLP-0596-S05
மேலும் நாம் நிறுத்துவோம்!{}; அதாவது
$A = g[\alpha]$ ஆகும் ஒரு (மீச்சிறிய) வரிசையெண் $\alpha$
உள்ளது. இல்லையென்றால், $g$ இன் நிறுத்துமுன் தொடக்கத் துண்டு
!!{injective} என்பதால், ஒவ்வொரு வரிசையெண் $\alpha$
க்கும் $\cardless{\alpha}{\Pow{A} \setminus \{\emptyset\}}$
கிடைக்கும். இது \olref[hartogs]{HartogsLemma} க்கு முரண்.
எனவே $\ran{g} = A$ என்பதும் நிறுத்துமுன் தொடக்கத்
துண்டுக்கு உண்மை.

இவற்றை ஒன்றுசேர்த்தால், ஏதேனும் ஒரு வரிசையெண்ணிலிருந்து
$A$ க்கு $g$ இன் நிறுத்துமுன் தொடக்கத் துண்டு
!!a{bijection} ஆகும். இந்த $g$ ஐப் பயன்படுத்தி $A$ ஐ நன்கு
வரிசைப்படுத்தலாம்.
\end{proof}

% SEGMENT OLP-0596-S06
ஆகவே நன்கு வரிசைப்படுத்தலும் தேர்வும் ஒன்றாகவே நிற்கும்
அல்லது வீழும். ஆனால் அவை நிற்குமா, வீழுமா என்ற கேள்வி
இன்னும் உள்ளது.

\end{document}
